% TOP_PROBLEM: {"problemId": "problem.combinatorial-boolean-matrix-multiplication-in-near-linear-time", "canonicalTitle": "Combinatorial Boolean Matrix Multiplication in Near-Linear Time", "releaseRank": 480, "primaryDomain": "theoretical_computer_science", "primaryDomainLabel": "Theoretical computer science", "displayStatus": "open", "releaseStatus": "open"}
% !TeX program = lualatex
% ATTEMPT_STATUS: unresolved
\documentclass[11pt]{article}
\usepackage[margin=25mm]{geometry}
\usepackage{fontspec}
\setmainfont{Latin Modern Roman}
\usepackage{amsmath,amssymb,mathtools}
\usepackage{unicode-math}
\setmathfont{Latin Modern Math}
\usepackage{xurl,hyperref}
\hypersetup{colorlinks=true,urlcolor=blue}
\setcounter{secnumdepth}{0}
\setlength{\parindent}{0pt}
\setlength{\parskip}{5pt}
\setlength{\emergencystretch}{3em}
\begin{document}
\section*{\texorpdfstring{Combinatorial Boolean Matrix Multiplication in Near-Linear Time}{Combinatorial Boolean Matrix Multiplication in Near-Linear Time}}\label{480}
\subsection{Definitions and mathematical statement}
For Boolean matrices $A,B\in\{0,1\}^{n\times n}$, their Boolean product is
\[
 C_{ij}=\bigvee_{k=1}^n(A_{ik}\wedge B_{kj}).
\]
The target asks for a purely combinatorial algorithm computing all $n^2$ output bits in $O(n^{2+o(1)})$ time. Here the exponent loss tends to zero as $n\to\infty$, so for every fixed ${\varepsilon}>0$ the bound is eventually $O(n^{2+{\varepsilon}})$.

The exclusion of algebraic Strassen-like tensor methods is part of the record. However, purely combinatorial is not itself a machine model or a universally formalized class of algorithms. A fully formal solution must specify allowed operations and justify that it meets the intended exclusion, using the cited computational convention. A generic fast algebraic matrix-multiplication algorithm is not silently counted as a solution to this restricted-method problem.
\par\smallskip\textit{Formulation and concept references:} \href{https://doi.org/10.1109/FOCS.2015.19}{[S1]}.

\subsection{Short English statement}
Can Boolean matrix multiplication be performed in nearly quadratic time using a genuinely combinatorial rather than algebraic tensor-based method?
\subsection{Research attempt}
\paragraph{Scope, computational convention, and source audit.}
This attempt by GPT-6 Astra Ultra was prepared on 27 September 2026.
The requested time is nearly linear in the $2n^2$ input bits, not in the
matrix dimension $n$. A bound $n^{3-o(1)}$ is still far from
$n^{2+o(1)}$. Yu's author manuscript [E1], Theorem 2, supplies a
combinatorial bound $\widehat O(n^3/\log^4 n)$, with logarithms of
logarithms suppressed. The newer primary manuscript [E2], Theorem 1.2,
gives a deterministic bound
\[
                       n^3/2^{\Omega((\log n)^{1/7})}.
\]
That saving exceeds every fixed power of $\log n$, but the resulting
exponent is still $3-o(1)$, not the exponent two requested here.
Section 1.1 of [E2] also explicitly discusses the absence of a universally
accepted formal definition of ``combinatorial''. No source inspected here
establishes the target. The publisher DOI in [S1] was not retrievable;
the linked author text [E1] was used to verify the actual theorem.

For the work below, use a deterministic word RAM with word size
$w=4\lceil\log_2(n+2)\rceil$. The table algorithm restricts its block
size to $b\leq\lfloor w/2\rfloor$, so its table of
$2^b\lceil n/w\rceil$ words, its pointers, and the input and output
matrices fit in the $2^w$-word address space. Fixed small dimensions may
be handled directly. This explicit restriction is needed: a mask fitting
in one word does not by itself imply that every vector indexed by that
mask can be stored in addressable memory.
Boolean operations, bounded shifts, comparisons, ordinary index arithmetic,
and array accesses on one word cost constant time. An $n$-bit vector uses
$W=\lceil n/w\rceil$ words. Reading an unpacked input or writing unpacked
output costs an additional $O(n^2)$, which is included. The algorithms
below use unions and intersections of represented sets, finite lookup
tables, and explicit enumeration; they invoke no algebraic matrix product
or tensor identity. This specifies the meaning of their stated bounds
without purporting to formalize every algorithm that the literature calls
combinatorial.

The attack has three parts: account fully for subset-table preprocessing,
separate output pairs from their potentially many witnesses, and investigate
whether sampling makes the remaining witness search inexpensive. The last
step produces a useful expectation bound but lacks the required reporting
algorithm. Its missing cost cannot be charged just to the output size.

\paragraph{Lemma 480.1: a fully accounted subset-table algorithm.}
For every integer $1\leq b\leq\min(n,\lfloor w/2\rfloor)$, Boolean multiplication can
be performed by the following combinatorial algorithm in
\begin{equation}\label{a480:tablecost}
 O\left(n^2+\left\lceil\frac nb\right\rceil
                         (2^b+n)\left\lceil\frac nw\right\rceil\right)
\end{equation}
word operations, using
$O(n^2/w+2^b\lceil n/w\rceil+n)$ words of working and matrix storage.

\emph{Algorithm and proof.} Partition the witness indices $[n]$ into
consecutive blocks $I$ of size at most $b$. For a fixed block, and each
subset $S\subseteq I$, store the $n$-bit row vector
\[
                     T_I(S)=\bigvee_{k\in S}B_{k,*}.
\]
Start with the zero vector at the empty set. For a nonempty mask $S$,
remove its least set bit $k$ and use
\begin{equation}\label{a480:recurrence}
                     T_I(S)=T_I(S\setminus\{k\})\vee B_{k,*}.
\end{equation}
The right side uses an already computed smaller mask. Each vector union
costs $O(W)$, so the table costs $O(2^{|I|}W)$ operations, including
initialization. Least-bit indices and table pointers can themselves be
prepared for all masks in $O(2^{|I|})$ time; an unsupported constant-time
operation on an arbitrarily large integer is not needed.

For each row $i$ of $A$, form its mask
$S_i=\{k\in I:A_{ik}=1\}$ and OR $T_I(S_i)$ into row $i$ of the
output. This costs $O(nW)$ operations per block. Masks use at most one
word; extraction from a packed row uses at most two adjacent words and
shifts, or all masks may be formed by a single $O(n^2)$ scan of unpacked
input. Repeat over all blocks. A bit $(i,j)$ is set exactly if some
block contains a witness $k$ with $A_{ik}=B_{kj}=1$, which is precisely
the Boolean product condition. This proves correctness.

There are $\lceil n/b\rceil$ blocks, giving
\eqref{a480:tablecost}. Tables can be discarded after every output row
has used the current block, so only one block's table occupies memory.
The last shorter block changes neither the recurrence nor the bound.
The output and the two input matrices use $O(n^2/w+n)$ words.
\hfill$\square$

Taking $b=\lfloor\log_2 n\rfloor$ for $n\geq2$, which satisfies
the stated mask and address-space restrictions, gives
\begin{equation}\label{a480:fourrussians}
                        O(n^2+n^3/\log^2 n).
\end{equation}
This rederives the classical Four Russians benchmark, not the stronger
frontier of [E2]. Both preprocessing and queries appear in the accounting.
Setting $b=n^{1-o(1)}$ in a symbolic query count
$n^3/(bw)$ would almost reach the desired scale, but would violate the
one-word mask convention and require $2^b$ stored vectors in this
implementation. Treating that table as a free oracle invalidates the
calculation. More generally, any polynomial-size table that stores all
$2^b$ masks forces $b=O(\log n)$.

\paragraph{An adaptive-table variant and its limitation.}
One might store only masks actually requested by rows of $A$. Let $q_I$
be their number in a block $I$. Close that family under repeatedly removing
the least set bit. The closure has at most $1+|I|q_I$ members: each of
the $q_I$ masks contributes a path of length at most $|I|$ to the empty
set. Computing the recurrence only on this closure is correct, with
\[
 O\left(W\sum_I\min\{2^{|I|},1+|I|q_I\}+nW\#\{I\}+n^2\right)
\]
time, provided masks and their dictionary operations are also charged.
For large masks, they require multiple words; the formula alone does not
erase that additional representation cost.

The hoped-for universal bound $q_I=n^{o(1)}$ is false. If $n=2^b$,
take the rows of $A$ restricted to a chosen block of $b$ columns to be
all $2^b$ masks. Then $q_I=n$. Repeating those bit patterns on many
blocks defeats that particular pattern-count bound simultaneously. This
does not prove a general lower bound: with $B$ the identity, the product
is just $A$, and a different algorithm should exploit that. It only shows
that fewer requested masks cannot be assumed without an additional
structural decomposition.

\paragraph{Lemma 480.2: an exact witness-sensitive algorithm.}
For $k\in[n]$, define
\[
 L_k=\{i:A_{ik}=1\},\qquad R_k=\{j:B_{kj}=1\},\qquad
 t_{ij}=|\{k:A_{ik}=B_{kj}=1\}|.
\]
Then the one-entries of $C$ form the union of the combinatorial rectangles
$L_k\times R_k$. In particular they can be computed in
\begin{equation}\label{a480:witnesscost}
 O(n^2+T)\quad\text{time},\qquad
 T=\sum_k|L_k||R_k|=\sum_{i,j}t_{ij}.
\end{equation}

\emph{Proof.} Build the incidence lists by scanning the input. For every
$k$, enumerate all pairs in $L_k\times R_k$ and set the corresponding
output bit. Every enumerated pair has a valid witness, and every positive
pair is enumerated at least once. The equality of the two expressions for
$T$ counts triples $(i,k,j)$ in two orders. Each update is constant time
even in a packed representation, and output initialization costs $O(n^2)$
or less.\hfill$\square$

Thus the requested scale is achieved on any family with
$T\leq n^{2+o(1)}$. This includes useful sparse configurations, but the
bound cannot be inferred from the $n^2$ output positions. If $A=B=J$
is the all-one matrix, $T=n^3$ although there are only $n^2$ positive
positions. After the first rectangle, every later rectangle is redundant;
the explicit pair-enumeration loop still executes its cubic number of
iterations. Recognizing $J\cdot J=J$ is easy, so this is a stress test
of the accounting rather than a hard instance for all algorithms.

The complementary issue occurs with rare witnesses. If $A=I$, every
positive output has exactly one witness and \eqref{a480:witnesscost} is
quadratic, even if $B$ is dense. Sampling a few witness indices, however,
misses most positive pairs. A correct combined approach needs to handle
both repetition and rare witnesses; neither density of $C$ nor density
of the individual input matrices determines which behavior occurs.

\paragraph{A rigorous special case using duplicate row and column types.}
Let $r$ be the number of distinct rows of $A$ and $c$ the number of
distinct columns of $B$. There is a deterministic algorithm with time
\begin{equation}\label{a480:types}
                 O(n^2+rc\lceil n/w\rceil).
\end{equation}
To see this, classify the $n$ input rows and columns by binary tries,
reading each of their $n$ bits once, for $O(n^2)$ total time and nodes.
For each representative pair, compute whether the bitwise intersection
is nonempty in $O(\lceil n/w\rceil)$ word operations. Its answer is
valid for every pair of indices in those two classes. Expand the table
of $rc$ answers to all $n^2$ positions. Storing the tries in an uncompressed
form costs $O(n^2)$ words, still polynomial; their construction is not
assumed free.

For example, $rc\leq nw\,n^{o(1)}$ gives the desired time scale for
this family. But both $r$ and $c$ can equal $n$. The all-one example has
$r=c=1$ and is easy for this algorithm, illustrating why the preceding
cubic enumeration example is not a complexity lower bound. A universal
algorithm needs a stronger notion of repeated structure than equality
of entire row or column vectors.

\paragraph{Sampling track: a quantitative residual bound with a missing operation.}
Temporarily allow independent uniform samples of witness indices, with
replacement. This is an exploratory randomized subroutine; it is not
presented as the deterministic solution. After $s\geq1$ samples, mark
the union of their rectangles. A pair with $t$ witnesses remains
unmarked with probability $(1-t/n)^s$. Therefore the expected number
of still relevant witness triples is
\begin{equation}\label{a480:residualmass}
 \mathbb E T_{\rm rem}
 =\sum_{i,j} t_{ij}(1-t_{ij}/n)^s
 \leq\frac{n^3}{s+1}.
\end{equation}
Indeed the function $x(1-x)^s$ on $[0,1]$ has maximum
$s^s/(s+1)^{s+1}\leq1/(s+1)$, found by differentiating or by
checking its single critical point. Substitute $x=t_{ij}/n$ and sum
over $n^2$ pairs. Pairs with no witness contribute zero, as they should.

A second useful bound is
\[
 \Pr[\text{some pair with }t_{ij}\geq t\text{ is missed}]
                         \leq n^2 e^{-st/n},
\]
by $(1-x)^s\leq e^{-sx}$ and a union bound. Thus high-witness pairs
are covered by sufficiently many samples. This assertion makes no promise
about the zero entries or about finding all the low-witness positive entries.

Marking one sampled rectangle by OR-ing its column bitset into its
designated rows costs at most $O(nW)$ operations. Suppose, as an
\emph{unproved extra hypothesis}, that after marking we could enumerate
only the residual triples in $O(n^2+T_{\rm rem})$ additional time.
Then the expected bound would be
\begin{equation}\label{a480:hypothetical}
                  O(n^2+sn^2/w+n^3/(s+1)).
\end{equation}
Even this hypothetical combination optimizes at
$s\asymp\sqrt{nw}$, giving $O(n^2+n^{5/2}/\sqrt w)$ rather than
nearly quadratic time. More importantly, the extra hypothesis is not
established. Iterating through every pair in every original rectangle and
testing a mark still costs $T$, not $T_{\rm rem}$. The first missing
algorithmic operation is reporting a rectangle's unmarked incidences
without scanning its already marked area. Bounding the number reported
does not bound the time to find them.

This is also why ``every output bit is set at most once'' is not a valid
runtime proof for a version that tests the same already-set bit many times.
A valid amortization needs to charge those unsuccessful tests, the data
structure updates, and any duplication of row or column lists caused by
recursion. The regularity and decomposition results in [E2] contain real
work of this kind; their efficiency is not automatically inherited by an
unspecified rectangle union routine.

\paragraph{Further stress points: certificates and the machine model.}
For each positive entry, supplying one witness $k$ gives a certificate
whose validity is checked by two input accesses. These certificates verify
that the claimed one-entries are genuine, but do not certify that a claimed
zero has no witness. Returning an arbitrary subset of the true one-entries
would pass that incomplete test. Producing and certifying the entire
product remain separate obligations.

Likewise, an $n$-bit intersection is not one operation on the specified
word RAM: it costs $\Theta(n/w)$ in the worst case. Taking $w=n$ would
give a different machine. Lookup tables must be charged when constructed
and stored, and a promise that their contents depend only on $n$ does
not license superpolynomial advice in a uniform algorithm. None of the
proofs above relies on that change of model. Their failure to reach the
target does not establish impossibility for the broader, informal class
of combinatorial algorithms.

\paragraph{Exact finite verification.}
The companion Python script compares the subset-table routine with direct
Boolean multiplication for every pair of matrices of sizes one, two, and
three, at every block size $1\leq b\leq n$. This is 262,404 matrix
pairs. It also checks the witness-count double sum on every pair, and
tests 80 further fixed-seed pairs in dimensions 4, 7, 16, and 31 with
several block sizes, including incomplete final blocks. The inequality
used in \eqref{a480:residualmass} is checked with exact rational arithmetic
for $1\leq n\leq29$, all integer witness counts, and $1\leq s\leq8$.

The implementation represents rows by Python integers to keep these small
functional checks short. Those integers are not being charged as
constant-time arbitrarily long words in the complexity proof: the proof
explicitly charges $W$ machine words per vector. The script is a
correctness diagnostic, not a performance benchmark or an empirical
asymptotic claim. Scripts and results are retained in
\nolinkurl{_Local/top_problems/continuation_488_477_2026_09_27/algorithms/}.

\paragraph{Outcome and next concrete obligations.}
\textbf{UNRESOLVED.} The completed work rederives the classical table
algorithm with full preprocessing and storage costs, proves two input-sensitive
algorithms, and quantifies what witness sampling leaves behind. It supplies
explicit examples invalidating free-table, one-write-per-output, and
few-distinct-pattern shortcuts. None of these results attains
$O(n^{2+o(1)})$ for arbitrary inputs, nor disproves that such a
combinatorial algorithm exists.

For the pursued route, the first gap is a reporting mechanism whose cost
depends on new rectangle incidences and controlled structural overhead,
rather than all repeated incidences. A second gap remains even with the
hypothetical reporting bound: the sampling tradeoff
\eqref{a480:hypothetical} is too large. Specific next targets are a
multiscale treatment of witness multiplicities, a proof bounding the
total duplicated data across rectangle splits, and a structural certificate
that detects dense easy pieces cheaply. Those bounds must be proved in
the stated word model before any near-quadratic conclusion is warranted.

\subsection{Sources}
\begin{itemize}
\item[S1]H. Yu, FOCS 2015: 109-119, 2015.
\url{https://doi.org/10.1109/FOCS.2015.19}.
\textit{Formulation source.}
\item[E1]Huacheng Yu, An Improved Combinatorial Algorithm for Boolean
Matrix Multiplication, arXiv:1505.06811; author manuscript dated
26 February 2017, Theorem 2 and Introduction.
\url{https://arxiv.org/abs/1505.06811}.
\url{https://www.cs.princeton.edu/~hy2/files/cbmm.pdf}.
\textit{Author text consulted 27 September 2026; the original DOI was not
retrievable in this session.}
\item[E2]Amir Abboud, Nick Fischer, Zander Kelley, Shachar Lovett, and
Raghu Meka, New Graph Decompositions and Combinatorial Boolean Matrix
Multiplication Algorithms, arXiv:2311.09095v2, 27 May 2024,
Theorem 1.2 and Sections 1.1 and 2.1.
\url{https://arxiv.org/html/2311.09095v2}.
\textit{Primary theorem and computational-convention source; consulted
27 September 2026.}

\end{itemize}
\end{document}
